\documentclass[11pt]{article}
\newcommand{\ListaTitulo}{Gabarito --- Lista 10}
% Preâmbulo compartilhado — MATD48, Listas de Exercícios 2026
% Usado por Lista{NN}.tex e Gabarito{NN}.tex via % Preâmbulo compartilhado — MATD48, Listas de Exercícios 2026
% Usado por Lista{NN}.tex e Gabarito{NN}.tex via % Preâmbulo compartilhado — MATD48, Listas de Exercícios 2026
% Usado por Lista{NN}.tex e Gabarito{NN}.tex via \input{preamble}

\usepackage[utf8]{inputenc}
\usepackage[T1]{fontenc}
\usepackage[brazil]{babel}
\usepackage{fullpage}
\usepackage{amsmath,amssymb}
\usepackage{enumitem}
\usepackage{xcolor}
\usepackage{fancyhdr}
\usepackage{titlesec}
\usepackage{listings}
\usepackage{hyperref}
\usepackage{booktabs}
\usepackage{graphicx}

\definecolor{matd48azul}{HTML}{0B4F6C}
\definecolor{matd48verde}{HTML}{2E8B57}

\titleformat{\section}{\normalfont\Large\bfseries\color{matd48azul}}{\thesection}{1em}{}
\titleformat{\subsection}{\normalfont\large\bfseries\color{matd48azul}}{\thesubsection}{1em}{}

\providecommand{\ListaTitulo}{Lista de Exercícios}

\setlength{\headheight}{14pt}
\pagestyle{fancy}
\fancyhf{}
\lhead{\small MATD48 --- Planejamento de Experimentos A}
\rhead{\small \ListaTitulo}
\cfoot{\thepage}
\renewcommand{\headrulewidth}{0.4pt}

\lstset{
  basicstyle=\ttfamily\small,
  breaklines=true,
  frame=single,
  columns=fullflexible,
  backgroundcolor=\color{gray!8},
  commentstyle=\color{matd48verde},
  keywordstyle=\color{matd48azul}\bfseries,
  language=R,
  extendedchars=true,
  literate=%
    {á}{{\'a}}1 {à}{{\`a}}1 {â}{{\^a}}1 {ã}{{\~a}}1
    {é}{{\'e}}1 {ê}{{\^e}}1 {í}{{\'i}}1 {ó}{{\'o}}1
    {ô}{{\^o}}1 {õ}{{\~o}}1 {ú}{{\'u}}1 {ç}{{\c c}}1
    {Á}{{\'A}}1 {À}{{\`A}}1 {Â}{{\^A}}1 {Ã}{{\~A}}1
    {É}{{\'E}}1 {Ê}{{\^E}}1 {Í}{{\'I}}1 {Ó}{{\'O}}1
    {Ô}{{\^O}}1 {Õ}{{\~O}}1 {Ú}{{\'U}}1 {Ç}{{\c C}}1
}

\hypersetup{colorlinks=true, linkcolor=matd48azul, urlcolor=matd48azul, citecolor=matd48azul}

\newcommand{\cabecalho}[2]{%
  \begin{center}
    {\Large\bfseries\color{matd48azul} #1}\\[0.3em]
    {\large #2}\\[0.3em]
    \rule{\linewidth}{0.6pt}
  \end{center}
  \vspace{0.5em}
}

\newenvironment{questao}[1]{\vspace{1em}\noindent\textbf{Questão #1.}\hspace{0.4em}}{\par}
\newenvironment{solucao}{\vspace{0.3em}\noindent\textit{Solução.}\hspace{0.4em}\color{black!85}}{\par\vspace{0.5em}}


\usepackage[utf8]{inputenc}
\usepackage[T1]{fontenc}
\usepackage[brazil]{babel}
\usepackage{fullpage}
\usepackage{amsmath,amssymb}
\usepackage{enumitem}
\usepackage{xcolor}
\usepackage{fancyhdr}
\usepackage{titlesec}
\usepackage{listings}
\usepackage{hyperref}
\usepackage{booktabs}
\usepackage{graphicx}

\definecolor{matd48azul}{HTML}{0B4F6C}
\definecolor{matd48verde}{HTML}{2E8B57}

\titleformat{\section}{\normalfont\Large\bfseries\color{matd48azul}}{\thesection}{1em}{}
\titleformat{\subsection}{\normalfont\large\bfseries\color{matd48azul}}{\thesubsection}{1em}{}

\providecommand{\ListaTitulo}{Lista de Exercícios}

\setlength{\headheight}{14pt}
\pagestyle{fancy}
\fancyhf{}
\lhead{\small MATD48 --- Planejamento de Experimentos A}
\rhead{\small \ListaTitulo}
\cfoot{\thepage}
\renewcommand{\headrulewidth}{0.4pt}

\lstset{
  basicstyle=\ttfamily\small,
  breaklines=true,
  frame=single,
  columns=fullflexible,
  backgroundcolor=\color{gray!8},
  commentstyle=\color{matd48verde},
  keywordstyle=\color{matd48azul}\bfseries,
  language=R,
  extendedchars=true,
  literate=%
    {á}{{\'a}}1 {à}{{\`a}}1 {â}{{\^a}}1 {ã}{{\~a}}1
    {é}{{\'e}}1 {ê}{{\^e}}1 {í}{{\'i}}1 {ó}{{\'o}}1
    {ô}{{\^o}}1 {õ}{{\~o}}1 {ú}{{\'u}}1 {ç}{{\c c}}1
    {Á}{{\'A}}1 {À}{{\`A}}1 {Â}{{\^A}}1 {Ã}{{\~A}}1
    {É}{{\'E}}1 {Ê}{{\^E}}1 {Í}{{\'I}}1 {Ó}{{\'O}}1
    {Ô}{{\^O}}1 {Õ}{{\~O}}1 {Ú}{{\'U}}1 {Ç}{{\c C}}1
}

\hypersetup{colorlinks=true, linkcolor=matd48azul, urlcolor=matd48azul, citecolor=matd48azul}

\newcommand{\cabecalho}[2]{%
  \begin{center}
    {\Large\bfseries\color{matd48azul} #1}\\[0.3em]
    {\large #2}\\[0.3em]
    \rule{\linewidth}{0.6pt}
  \end{center}
  \vspace{0.5em}
}

\newenvironment{questao}[1]{\vspace{1em}\noindent\textbf{Questão #1.}\hspace{0.4em}}{\par}
\newenvironment{solucao}{\vspace{0.3em}\noindent\textit{Solução.}\hspace{0.4em}\color{black!85}}{\par\vspace{0.5em}}


\usepackage[utf8]{inputenc}
\usepackage[T1]{fontenc}
\usepackage[brazil]{babel}
\usepackage{fullpage}
\usepackage{amsmath,amssymb}
\usepackage{enumitem}
\usepackage{xcolor}
\usepackage{fancyhdr}
\usepackage{titlesec}
\usepackage{listings}
\usepackage{hyperref}
\usepackage{booktabs}
\usepackage{graphicx}

\definecolor{matd48azul}{HTML}{0B4F6C}
\definecolor{matd48verde}{HTML}{2E8B57}

\titleformat{\section}{\normalfont\Large\bfseries\color{matd48azul}}{\thesection}{1em}{}
\titleformat{\subsection}{\normalfont\large\bfseries\color{matd48azul}}{\thesubsection}{1em}{}

\providecommand{\ListaTitulo}{Lista de Exercícios}

\setlength{\headheight}{14pt}
\pagestyle{fancy}
\fancyhf{}
\lhead{\small MATD48 --- Planejamento de Experimentos A}
\rhead{\small \ListaTitulo}
\cfoot{\thepage}
\renewcommand{\headrulewidth}{0.4pt}

\lstset{
  basicstyle=\ttfamily\small,
  breaklines=true,
  frame=single,
  columns=fullflexible,
  backgroundcolor=\color{gray!8},
  commentstyle=\color{matd48verde},
  keywordstyle=\color{matd48azul}\bfseries,
  language=R,
  extendedchars=true,
  literate=%
    {á}{{\'a}}1 {à}{{\`a}}1 {â}{{\^a}}1 {ã}{{\~a}}1
    {é}{{\'e}}1 {ê}{{\^e}}1 {í}{{\'i}}1 {ó}{{\'o}}1
    {ô}{{\^o}}1 {õ}{{\~o}}1 {ú}{{\'u}}1 {ç}{{\c c}}1
    {Á}{{\'A}}1 {À}{{\`A}}1 {Â}{{\^A}}1 {Ã}{{\~A}}1
    {É}{{\'E}}1 {Ê}{{\^E}}1 {Í}{{\'I}}1 {Ó}{{\'O}}1
    {Ô}{{\^O}}1 {Õ}{{\~O}}1 {Ú}{{\'U}}1 {Ç}{{\c C}}1
}

\hypersetup{colorlinks=true, linkcolor=matd48azul, urlcolor=matd48azul, citecolor=matd48azul}

\newcommand{\cabecalho}[2]{%
  \begin{center}
    {\Large\bfseries\color{matd48azul} #1}\\[0.3em]
    {\large #2}\\[0.3em]
    \rule{\linewidth}{0.6pt}
  \end{center}
  \vspace{0.5em}
}

\newenvironment{questao}[1]{\vspace{1em}\noindent\textbf{Questão #1.}\hspace{0.4em}}{\par}
\newenvironment{solucao}{\vspace{0.3em}\noindent\textit{Solução.}\hspace{0.4em}\color{black!85}}{\par\vspace{0.5em}}


\begin{document}

\cabecalho{Gabarito --- Lista de Exercícios 10}{Teste de Friedman e blocos incompletos balanceados (Capítulo 4 / Aula 10)}

\begin{questao}{1} \textbf{(Psicologia --- estatística de Friedman)}
\end{questao}
\begin{solucao}
\begin{enumerate}[label=(\alph*)]
  \item Cada um dos $b=10$ estudantes distribui os postos $1,2,3$ (uma vez cada) entre as três
    técnicas, então a soma de \emph{todos} os postos atribuídos, somando as 3 técnicas, é
    $b\cdot(1+2+3) = 10\times 6 = 60$. De fato, $R_1+R_2+R_3=12+20+28=60$. \checkmark
  \item
    $$
    F_r = \frac{12}{bt(t+1)}\sum_i R_i^2 - 3b(t+1)
        = \frac{12}{10\cdot3\cdot4}(12^2+20^2+28^2) - 3\cdot10\cdot4
        = \frac{12}{120}(144+400+784) - 120
    $$
    $$
    = 0{,}1\times1328 - 120 = 132{,}8-120 = \textbf{12{,}8}.
    $$
  \item Como $F_r=12{,}8 > \chi^2_{0{,}05;2}=5{,}99$, \textbf{rejeita-se $H_0$}: há evidência de
    diferença sistemática entre as três técnicas de estudo (a técnica associada a $R_3=28$ tende
    a ser mais frequentemente a "melhor" dentro de cada estudante).
\end{enumerate}
\end{solucao}

\begin{questao}{2} \textbf{(Dedução --- valor esperado da soma de postos sob $H_0$)}
\end{questao}
\begin{solucao}
\begin{enumerate}[label=(\alph*)]
  \item Sob $H_0$, as $t!$ ordenações possíveis dos $t$ tratamentos dentro de um bloco são
    igualmente prováveis (nenhum tratamento é sistematicamente favorecido). Fixado um tratamento
    $i$, cada um dos $t$ postos possíveis $1,\dots,t$ aparece associado a $i$ em exatamente
    $(t-1)!$ das $t!$ ordenações (basta permutar os $t-1$ tratamentos restantes nas demais
    posições), a mesma contagem para todo posto. Logo, marginalmente, o posto de $i$ é uniforme
    sobre $\{1,\dots,t\}$.
  \item Pelo item (a), $\text{posto}(Y_{ij}) \sim \text{Uniforme}\{1,\dots,t\}$, logo
    $$
    \mathbb E[\text{posto}(Y_{ij})] = \frac{1}{t}\sum_{k=1}^t k = \frac{1}{t}\cdot\frac{t(t+1)}{2}
    = \frac{t+1}{2}.
    $$
  \item Como os $b$ blocos são independentes, $R_i=\sum_{j=1}^b \text{posto}(Y_{ij})$ é soma de
    $b$ variáveis com a mesma esperança $\frac{t+1}{2}$, logo
    $$
    \mathbb E[R_i] = b\cdot\frac{t+1}{2}.
    $$
    Com $b=10$, $t=3$: $\mathbb E[R_i]=10\times2=\textbf{20}$ para cada técnica, sob $H_0$. Os
    valores observados $R=(12,20,28)$ do Exercício 1 se afastam bastante e de forma assimétrica
    dessa expectativa comum de 20 — consistente com a rejeição de $H_0$ obtida ali.
\end{enumerate}
\end{solucao}

\begin{questao}{3} \textbf{(Agricultura --- parâmetros de um BIB)}
\end{questao}
\begin{solucao}
\begin{enumerate}[label=(\alph*)]
  \item De $\lambda(t-1)=r(k-1)$ com $t=6,k=3$: $5\lambda = 2r \Rightarrow r = \dfrac{5\lambda}{2}$.
    Para $r$ inteiro, $\lambda$ deve ser par; o menor valor positivo é $\lambda=2$, dando
    $r=\dfrac{5\times2}{2}=5$. De $bk=tr$: $b=\dfrac{tr}{k}=\dfrac{6\times5}{3}=10$.
    Portanto $(t,b,r,k,\lambda)=(6,10,5,3,2)$.
  \item Verificação: $bk=10\times3=30$ e $tr=6\times5=30$ \checkmark;
    $\lambda(t-1)=2\times5=10$ e $r(k-1)=5\times2=10$ \checkmark. Ambas as identidades se
    confirmam.
\end{enumerate}
\end{solucao}

\begin{questao}{4} \textbf{(Agricultura --- verificando e completando um BIB)}
\end{questao}
\begin{solucao}
\begin{enumerate}[label=(\alph*)]
  \item Contando as ocorrências: variedade 1 aparece nos blocos 1, 2 e 3 ($r_1=3$); variedade 2
    aparece nos blocos 1 e 2 ($r_2=2$); variedade 3 aparece nos blocos 1 e 3 ($r_3=2$);
    variedade 4 aparece nos blocos 2 e 3 ($r_4=2$). Como $r_1\ne r_2=r_3=r_4$, \textbf{este
    arranjo de 3 blocos não é um BIB}. A condição de $r$ constante é necessária porque, em um
    BIB, todo tratamento deve ter a mesma precisão (o mesmo número de repetições); sem isso, as
    variâncias das médias estimadas de tratamento já diferem antes mesmo de considerar os pares —
    por isso a condição é necessária, mas não suficiente: mesmo com $r$ igual para todos, ainda
    seria preciso conferir que todo \emph{par} de tratamentos coocorre o mesmo número de vezes.
  \item Basta acrescentar o bloco que contém exatamente as três variedades ainda não combinadas
    entre si na proporção correta: $\text{Bloco 4}=\{2,3,4\}$. Com os 4 blocos
    $\{1,2,3\},\{1,2,4\},\{1,3,4\},\{2,3,4\}$, cada variedade aparece em $r=3$ blocos, e cada par
    coocorre em exatamente $\lambda=2$ blocos (por exemplo, o par $(1,2)$ aparece nos blocos 1 e
    2; o par $(2,3)$ aparece nos blocos 1 e 4; e assim por diante para os $\binom{4}{2}=6$ pares).
    Este é exatamente o BIB $(t,b,r,k,\lambda)=(4,4,3,3,2)$ apresentado em aula.
\end{enumerate}
\end{solucao}

\begin{questao}{5} \textbf{(Dissertativa --- Friedman vs. DBCA, blocos completos vs. incompletos)}
\end{questao}
\begin{solucao}
\begin{enumerate}[label=(\alph*)]
  \item Quando a normalidade de fato vale, o teste de Friedman é \textbf{menos sensível} (tem
    menor poder) que a ANOVA do DBCA, porque descarta a magnitude das diferenças entre respostas
    e usa apenas sua ordem dentro de cada bloco — informação estritamente menor do que os valores
    numéricos completos. Com $b=20$ blocos e dados bem comportados, o pesquisador provavelmente
    não muda a decisão qualitativa (rejeitar ou não $H_0$), mas paga o preço de um teste menos
    poderoso do que precisaria ser: mais provável não detectar uma diferença real de magnitude
    moderada (erro tipo II mais frequente) do que se tivesse usado a ANOVA apropriada para dados
    normais.
  \item Blocos incompletos balanceados fazem sentido mesmo quando o tamanho do bloco não é uma
    restrição física, mas quando incluir todos os $t$ tratamentos em cada bloco comprometeria a
    \textbf{homogeneidade interna do bloco} que justifica blocar em primeiro lugar — por exemplo,
    se blocos maiores (para caber mais tratamentos) precisassem se espalhar por uma área com mais
    heterogeneidade de solo, ou se a sessão de avaliação ficasse tão longa que introduzisse fadiga
    ou efeitos de ordem não controlados. Nesses casos, blocos incompletos, porém pequenos e
    homogêneos, podem produzir comparações mais precisas do que blocos completos, porém
    heterogêneos.
\end{enumerate}
\end{solucao}

\end{document}
